% =============================================================================
\section{Comparative Analysis}
\label{sec:comparative}
% =============================================================================

This section provides a systematic comparison of the PURECORTEX tokenomic
framework against the prevailing agent tokenization architecture, referred to
as ``Platform~A'' (see Section~\ref{subsec:rw-agent}). We evaluate 15
dimensions spanning allocation philosophy, creator incentives, revenue
distribution, economic mechanics, governance, and technical infrastructure.

\subsection{Comparison Framework}
\label{subsec:comp-framework}

Table~\ref{tab:comparison} presents the full 14-dimension comparison.

\begin{table}[H]
\centering
\caption{15-dimension comparative analysis: PURECORTEX vs.\ Platform~A.}
\label{tab:comparison}
\small
\begin{tabular}{@{}p{2.8cm}p{4.8cm}p{4.8cm}@{}}
\toprule
\textbf{Dimension} & \textbf{PURECORTEX} & \textbf{Platform A} \\
\midrule
Insider Allocation & 10\% creator only (0\% VC, 0\% team/advisors) & Significant team + VC allocation \\
\addlinespace
Creator Vesting & 10\% at TGE + 90\% daily over 180 days & 100\% immediate, no vesting \\
\addlinespace
Revenue to Holders & 90\% of ALL revenue $\to$ buyback-burn (Assistance Fund) & 2-way split (creator + protocol) \\
\addlinespace
Staking Timing & Available from day one (launch parity) & Introduced post-launch; early adopters disadvantaged \\
\addlinespace
Fee Model & Dynamic: 0.5\%--2.0\% (volatility-adaptive) & Flat 1\% on all transactions \\
\addlinespace
LP Lock & 10-year LogicSig (immutable, no admin keys) & LP locked (mechanism varies) \\
\addlinespace
Supply Schedule & 4-year declining emissions + continuous buyback-burn & Fixed supply, no emission or burn schedule \\
\addlinespace
Burn Mechanics & 90\% of all revenue $\to$ Assistance Fund $\to$ continuous burn & No systematic burn \\
\addlinespace
Composability & MCP tool protocol + x402 micropayments + CS rewards & No inter-agent composability \\
\addlinespace
Governance Model & On-chain: Senator AI + veCORTEX lawmakers + constitution & Centralized team control \\
\addlinespace
Constitutional Framework & Immutable preamble + 7 amendable articles, on-chain & None \\
\addlinespace
Security Model & Oracle-free, atomic txns, no admin keys & EVM-based (oracle + reentrancy surface) \\
\addlinespace
Chain Finality & Instant ($<3.3$s, deterministic) & Probabilistic (block confirmations) \\
\addlinespace
Agent Interop & MCP standard + composability rewards & Isolated agents, no interop \\
\addlinespace
Oracle Dependency & None (all inputs on-chain) & External price feeds for some operations \\
\bottomrule
\end{tabular}
\end{table}

\subsection{Detailed Analysis by Dimension}
\label{subsec:dim-analysis}

\subsubsection{Creator Vesting}

The absence of creator vesting in Platform~A creates a well-documented
moral hazard: creators can liquidate their entire allocation immediately after
launch, extracting value from buyers who expect continued development. This
has been empirically associated with the ``pump-and-dump'' phenomenon in agent
token markets, wherein agent tokens experience rapid price appreciation
followed by creator exit and price collapse.

PURECORTEX's 180-day vesting schedule addresses this by ensuring that 99\% of
the creator's allocation remains locked for six months. The creator's optimal
strategy under vesting is to continue developing the agent (which supports
token price and trading volume, generating fee revenue), rather than abandoning
it (which causes token price decline but cannot be monetized due to the lock).
This alignment is formalized in Proposition~\ref{prop:creator-wealth} and
Theorem~\ref{thm:separating}.

\subsubsection{Staking Timing}

When staking is introduced after token launch, the earliest participants---who
bore the maximum price-discovery risk---receive no preferential staking
treatment. Late entrants who waited for price stability can lock at the same
terms. This temporal inequity discourages early participation, reducing initial
liquidity and price discovery quality.

PURECORTEX's day-one staking availability ensures that early participants can
immediately lock their tokens, capturing the highest yield (year-1 emission
weight of 0.4) and the exclusive-period revenue (Proposition~\ref{prop:day-one}).
This creates a strong incentive for early capital commitment.

\subsubsection{Fee Model}

A flat 1\% fee is suboptimal in two regimes:
\begin{itemize}
  \item In low-volatility markets, 1\% is unnecessarily high, suppressing
    trading volume.
  \item In high-volatility markets, 1\% is insufficient to compensate for
    adverse selection, causing liquidity providers to withdraw.
\end{itemize}

The dynamic fee (Section~\ref{sec:dynamic-fees}) adapts to both regimes,
generating 23--31\% more cumulative revenue
(Section~\ref{subsec:sim-fees}) while maintaining lower friction during calm
periods.

\subsubsection{Revenue to Holders}

Platform~A distributes revenue through a simple two-way split between the
creator and the protocol, with the protocol's share controlled by the
centralized team. Token holders have no direct claim on protocol revenue and
benefit only through speculative price appreciation.

PURECORTEX directs 90\% of \emph{all} protocol revenue to the Assistance Fund,
which continuously buys CORTEX from the open market and burns it. This
mechanism creates a direct, automated, and transparent link between protocol
activity and token holder value: every fee generated reduces the circulating
supply, increasing per-token value for all holders. The model is inspired by
the approach of the leading decentralized perpetuals protocol, which
demonstrated that directing the vast majority of revenue to buyback-burn
produces stronger holder alignment than dividend-style distributions.

\subsubsection{Burn Mechanics}

Platform~A operates with a fixed supply and no systematic burn mechanism. Over
time, lost tokens (sent to incorrect addresses, abandoned wallets) represent
the only supply reduction---an unpredictable and uncontrolled process.

PURECORTEX implements continuous burns through the Assistance Fund
(Section~\ref{sec:buyback-burn}), with 90\% of all revenue funding buyback-burn
operations. This creates an aggressive deflationary trajectory with a crossover
at month 12--18 under moderate growth (Theorem~\ref{thm:deflationary},
Section~\ref{subsec:sim-supply}).

\subsubsection{Composability}

Platform~A agents operate in isolation. There is no standardized protocol for
one agent to invoke another's capabilities, no economic incentive for
inter-agent integration, and no mechanism for measuring or rewarding
composability. This architectural choice limits the ecosystem to a collection
of independent agents, forgoing the superlinear network effects predicted by
Metcalfe's law.

PURECORTEX's MCP-based composability (Section~\ref{sec:composability}) creates
an agent network where each new agent increases the value of the entire
ecosystem. The Composability Score and reward mechanism incentivize tool
exposure and integration, with the anti-sybil properties of
Proposition~\ref{prop:sybil} preventing gaming.

\subsubsection{Governance}

Centralized governance creates platform risk: users have no recourse if the
development team makes decisions that benefit the team at the expense of users.
There is no transparency requirement, no proposal process, and no
constitutional framework constraining team decisions.

PURECORTEX's on-chain governance (Section~\ref{sec:governance}) provides
transparent, constitutional governance with formal separation of powers. The
Senator AI ensures informed proposal quality; the Lawmaker system ensures
democratic legitimacy; the constitutional framework ensures rights protection.

\subsubsection{Chain Properties}

Algorand's instant finality eliminates confirmation-time uncertainties that
affect EVM-based platforms. An agent token purchase on PURECORTEX is
irreversible within 3.3 seconds; on platforms using probabilistic finality,
the same transaction may require 12--60 seconds of confirmation time before
the token balance can be used in subsequent operations (composability calls,
staking, governance voting).

The ASA model provides native asset support without the overhead of ERC-20
contract deployment, and atomic transaction groups enable the graduation
mechanism without reentrancy risk.

\subsection{Quantitative Impact Summary}
\label{subsec:quant-impact}

Table~\ref{tab:impact-summary} summarizes the quantitative impact of each
PURECORTEX design choice relative to the Platform~A baseline, as measured by
the Monte Carlo simulations in Section~\ref{sec:simulations}.

\begin{table}[H]
\centering
\caption{Quantitative impact of PURECORTEX design choices.}
\label{tab:impact-summary}
\begin{tabular}{@{}lrl@{}}
\toprule
\textbf{Design Choice} & \textbf{Impact} & \textbf{Metric} \\
\midrule
Creator vesting & $-67\%$ & Creator exit rate (180-day) \\
0\% insider allocation & $90\%$ & Community-controlled supply at TGE \\
Day-one staking & $+3\text{--}5\times$ & Early staker lifetime yield \\
Dynamic fees & $+27\%$ & Cumulative protocol revenue \\
90\% revenue buyback-burn & $-25\%$ & Circulating supply at month 48 \\
Composability rewards & $+8\times$ & Inter-agent connection density \\
On-chain governance & $0.38$ & Voting power Gini (vs.\ 1.0 centralized) \\
Crash recovery & $-43\%$ & Time to pre-crash TVL recovery \\
\bottomrule
\end{tabular}
\end{table}
